\section{Definición}
\label{sec:definicion}

Los autómatas celulares, concebidos por John von Neumann en 1948~\cite{damerow2010_neumann}, son una colección de células donde cada una de ellas puede adquirir un estado y se agrupan en una red. Además, estas células evolucionan en intervalos y cada una de ellas actualiza su estado en función del estado de sus células vecinas~\cite{weisstein2002_ca,wolfram1982_ca,fsancho2016_ca}.

Siguiendo de manera general a~\cite{bhattacharjee2019_ca,sayama2024_ca,Ilachinski2001_ca}, y considerando una red regular de células, un conjunto finito de estados y evoluciones en pasos discretos de tiempo, se define a un autómata celular como:

\begin{definition}
\label{def:ca}
Un autómata celular es una 4-tupla $(\mathcal{L}, \mathcal{S}, \mathcal{N}, \mathcal{R})$ donde:
\begin{itemize}
	\item $\mathcal{L} \subseteq \mathbb{Z}^D$ es el espacio celular $D$-dimensional. El conjunto de células vive en el espacio $\mathcal{L}$ y cada célula $\vec{i}$ se define como un vector $D$-dimensional que representa su posición en el espacio.

	\item $\mathcal{S}$ es un conjunto de $k$ estados. Cada célula $\vec{i}$ solo puede adquirir un valor de un conjunto finito de diferentes estados:
	\[
	\sigma_{\vec{i}}(t) \in \mathcal{S} \equiv \{ s_0, s_1, s_2, \dots, s_{k-1} \}, \; \forall \vec{i} \in \mathcal{L}
	\]
	donde $\sigma_{\vec{i}}(t)$ es el estado de la célula $\vec{i}$ en el tiempo $t$ y $\mathcal{S}$ es un anillo conmutativo finito, que generalmente se define por los enteros módulo $k$ ($\mathbb{Z}_k$).

	\item $\mathcal{N} = (\vec{v}_1, \vec{v}_2, \dots, \vec{v}_m)$ es el conjunto de los $m$ vectores de desplazamiento que definen la vecindad de una célula en $\mathcal{L}$. Así, los vecinos de la célula ubicada en $\vec{i} \in \mathcal{L}$ se encuentran en las posiciones $\vec{i} + \vec{v}, \; \forall \vec{v} \in \mathcal{N}$.

	\item $\mathcal{R} : \mathcal{S}^m \to \mathcal{S}$ es la regla local del autómata:
	\[
	\sigma_{\vec{i}}(t + 1) = \mathcal{R}\left( \{ \sigma_{\vec{i} + \vec{v}_j}(t) \}_{j=1}^{m} \right)
	\]
	que asigna el nuevo estado de una célula a partir de los estados de sus $m$ vecinos.
\end{itemize}
\end{definition}

\subsection{Vecindad}

La vecindad está estrechamente relacionada con la dimensionalidad del autómata celular. John von Neumann propuso para $\mathbb{Z}^2$ una vecindad de cinco células, que considera a la célula y a sus cuatro vecinas ortogonales más cercanas~\cite{neumann1966_ca}:
\[
\mathcal{N}_{\text{VN}} = ((0,-1), (-1,0), (0,0), (1,0), (0,1)).
\]
Moore añade los cuatro vecinos diagonales, dando una vecindad de nueve elementos~\cite{Kwon2015_moore}, utilizada por John Conway en el Juego de la Vida~\cite{lifewiki_life}. En $\mathbb{Z}^1$, la vecindad de radio $r$ contiene a todas las células a distancia menor o igual a $r$, incluyendo a la propia célula ($|\mathcal{N}| = 2r + 1$); con $r = 1$ fue utilizada por Stephen Wolfram para definir los autómatas celulares elementales~\cite{Wolfram2002_nks}.

\subsection{Condiciones de frontera}

Idealmente, los autómatas celulares evolucionan en espacios infinitos; sin embargo, se pueden establecer fronteras que restringen el espacio celular $\mathcal{L}$~\cite{bhattacharjee2019_ca,ponce2015_ca}.

\begin{definition}
Un autómata celular es finito si el espacio celular $\mathcal{L} \subseteq \mathbb{Z}^D$ es finito; en caso contrario, se denomina infinito.
\end{definition}

Un autómata celular finito siempre tiene fronteras en los límites de su espacio. Su comportamiento se conoce como condición de frontera, la cual suele ser abierta (se asigna un estado fijo a los vecinos faltantes) o periódica (los vecinos faltantes de una célula en la frontera son las células de la frontera opuesta)~\cite{ponce2015_ca}. Un autómata infinito con una configuración periódica equivale a un autómata finito con frontera periódica~\cite{bhattacharjee2019_ca}.

\begin{definition}
Sea $\mathcal{C} = \mathcal{S}^{\mathcal{L}}$ el conjunto de todas las configuraciones. Una configuración $\vec{\sigma} \in \mathcal{C}$ de un autómata celular $D$-dimensional es periódica si es invariante ante un conjunto de traslaciones linealmente independientes en $D$ direcciones:
\[
\vec{\sigma} = T_i^{n_i}(\vec{\sigma}), \; \forall i \in \{1, 2, \dots, D\}
\]
donde $T_i^{n_i}$ denota un desplazamiento de $n_i$ posiciones en la dimensión $i$.
\end{definition}

\subsection{Regla local}

Para autómatas celulares unidimensionales con radio $r$, la regla local se escribe como:
\[
\sigma_i(t + 1) = \mathcal{R}(\sigma_{i-r}(t), \dots, \sigma_i(t), \dots, \sigma_{i+r}(t)), \quad \mathcal{R} : \mathcal{S}^{2r+1} \to \mathcal{S}.
\]

Un ejemplo es el Juego de la Vida~\cite{lifewiki_life}, con $\mathcal{S} = \{0, 1\}$ (células muertas y vivas) y la vecindad de Moore, por lo que $\mathcal{R} : \{0,1\}^9 \to \{0,1\}$. Una célula muerta nace si tiene exactamente tres vecinos vivos; una célula viva sobrevive si tiene dos o tres vecinos vivos, y en cualquier otro caso muere. Esta regla es semitotalística, ya que depende de la suma de los estados de la vecindad, pero distingue a la célula central del resto~\cite{genaro2022_semitotalistic}.

Una regla es totalística si depende únicamente de la suma de los estados de todos los vecinos, sin distinguir posiciones~\cite{weisstein2002_totalistic}; hay $(|\mathcal{S}| - 1)|\mathcal{N}| + 1$ posibles valores de suma. Una regla es arbitraria si distingue la posición de cada elemento de la vecindad. Así:
\[
\mathcal{R}_{\text{totalísticas}} \subset \mathcal{R}_{\text{semitotalísticas}} \subset \mathcal{R}_{\text{arbitrarias}}.
\]

El autómata celular elemental considera $k = 2$, radio $r = 1$ ($|\mathcal{N}| = 3$) y la regla $\mathcal{R} : \{0,1\}^3 \to \{0,1\}$. Existen $k^{|\mathcal{N}|} = 2^3 = 8$ patrones posibles en la vecindad, y a cada uno se le asigna un nuevo estado, por lo que el número total de reglas es $k^{k^{|\mathcal{N}|}} = 2^8 = 256$. Cada regla se identifica con un número decimal leyendo como número binario los estados asignados a los patrones $111, 110, \dots, 000$; por ejemplo, la regla $30$ corresponde a $00011110$.

Nótese que la regla decimal solo es única cuando $|\mathcal{N}|$ y $|\mathcal{S}|$ son fijos: dos reglas que difieren en vecindad o en número de estados pueden compartir el mismo número. Por ejemplo, la regla $10$ con $|\mathcal{N}| = 2$ y $|\mathcal{S}| = 2$ es $1010$, mientras que con $|\mathcal{N}| = 3$ y $|\mathcal{S}| = 2$ es $00001010$, una regla completamente diferente.
